\chapter{Radial Quantization and the State-Operator Correspondence}
\label{ch:volume-iii-radial-quantization}

\section{The Weyl Map to the Cylinder}

Let \(x^\mu\) be Euclidean coordinates on \(\R^D\), and write
\[
  r=|x|,\qquad
  x^\mu=r n^\mu,\qquad
  n^\mu n_\mu=1.
\]
The flat metric is
\[
  \dd s^2
  =
  \dd r^2+r^2\dd\Omega_{D-1}^2,
\]
where \(\dd\Omega_{D-1}^2\) is the unit metric on \(S^{D-1}\).  Define
\[
  \tau=\log r .
\]
Then
\begin{equation}
  \frac{\dd s^2}{r^2}
  =
  \dd\tau^2+\dd\Omega_{D-1}^2.
  \label{eq:radial-cylinder-metric}
\end{equation}
Thus the punctured space \(\R^D\setminus\{0\}\) is Weyl equivalent to the
Euclidean cylinder
\[
  \R_\tau\times S^{D-1}.
\]
The origin corresponds to \(\tau=-\infty\), and infinity corresponds to
\(\tau=+\infty\).

\begin{figure}[H]
  \centering
  \resizebox{0.94\textwidth}{!}{%
  \begin{tikzpicture}[x=1cm,y=1cm,every node/.style={font=\scriptsize}]
    \begin{scope}
      \node[font=\small] at (0,2.35) {flat space};
      \draw[thick] (-2.2,-1.3) rectangle (2.2,1.5);
      \draw[blue!60!black,thick] (0,0) circle (0.55);
      \draw[blue!60!black,thick] (0,0) circle (1.05);
      \draw[blue!60!black,thick] (0,0) circle (1.55);
      \draw[orange!85!black,thick,-{Latex[length=2mm]}] (0,0)--(1.2,0.9);
      \node[orange!85!black] at (0.78,0.70) {\(r\)};
      \fill[purple!75!black] (0,0) circle (2pt);
      \node[purple!75!black,below left] at (0,0) {\(0\)};
    \end{scope}

    \draw[-{Latex[length=2.5mm]},thick] (2.75,0.12)--(4.05,0.12);
    \node[align=center] at (3.40,0.63) {\(\tau=\log r\)\\Weyl rescale};

    \begin{scope}[shift={(6.3,0)}]
      \node[font=\small] at (0,2.35) {cylinder};
      \draw[thick] (-1.35,-1.45)--(-1.35,1.55);
      \draw[thick] (1.35,-1.45)--(1.35,1.55);
      \draw[thick] (-1.35,1.55) arc[start angle=180,end angle=360,x radius=1.35,y radius=0.28];
      \draw[thick,dashed] (-1.35,1.55) arc[start angle=180,end angle=0,x radius=1.35,y radius=0.28];
      \draw[thick] (-1.35,-1.45) arc[start angle=180,end angle=360,x radius=1.35,y radius=0.28];
      \draw[thick,dashed] (-1.35,-1.45) arc[start angle=180,end angle=0,x radius=1.35,y radius=0.28];
      \foreach \y in {-0.8,0.0,0.8} {
        \draw[blue!60!black,thick] (-1.35,\y) arc[start angle=180,end angle=360,x radius=1.35,y radius=0.28];
        \draw[blue!60!black,thick,dashed] (-1.35,\y) arc[start angle=180,end angle=0,x radius=1.35,y radius=0.28];
      }
      \draw[-{Latex[length=2mm]},orange!85!black,thick] (-1.75,-1.35)--(-1.75,1.45);
      \node[orange!85!black,left] at (-1.75,0.15) {\(\tau\)};
      \node at (0,-1.95) {\(S^{D-1}\)};
    \end{scope}
  \end{tikzpicture}
  }
  \caption{Radial quantization identifies dilatations on \(\R^D\) with
  translations in cylinder time \(\tau=\log r\).}
  \label{fig:radial-cylinder-map}
\end{figure}

\section{States from Local Operators}

Let \(\mathcal O(0)\) be a local operator inserted at the origin.  In the
Euclidean path integral on a ball of radius \(R\), the insertion defines a
state on the boundary sphere \(S_R^{D-1}\):
\[
  \ket{\mathcal O;R}
  =
  \int_{\Phi|_{S_R}=\varphi}
  [D\Phi]_{B_R}\,
  \ee^{-S[\Phi]}\,
  \mathcal O(0).
\]
The right side is a wavefunctional of the boundary field configuration
\(\varphi\).  After the Weyl map to the cylinder, changing \(R\) translates
the state in \(\tau\).  The state associated with \(\mathcal O\) is denoted
\[
  \ket{\mathcal O}.
\]
Conversely, a state on \(S^{D-1}\) evolved to \(\tau=-\infty\) determines a
local insertion at the origin when the limit exists within the local operator
space.  This gives the state-operator map
\begin{equation}
  \mathcal V_{\mathrm{loc}}
  \longleftrightarrow
  \mathcal H_{S^{D-1}},
  \qquad
  \mathcal O\longmapsto\ket{\mathcal O}.
  \label{eq:state-operator-map}
\end{equation}

The vacuum state corresponds to the identity operator:
\[
  \mathbf 1\longleftrightarrow\ket{\Omega}.
\]
Products of local operators inserted at the origin are defined by the
short-distance operator product expansion and correspond to the corresponding
states after radial ordering.

\section{Dilatation as the Cylinder Hamiltonian}

The dilatation vector field on \(\R^D\) is
\[
  x^\mu\partial_\mu=r\partial_r=\partial_\tau.
\]
Therefore the dilatation charge maps to the Hamiltonian generating
translations in cylinder time.  With Lorentzian cylinder time
\[
  t=-\ii\tau,
\]
write the cylinder Hamiltonian as \(\widehat H_{\mathrm{cyl}}\).  The
corresponding operator relation is
\begin{equation}
  \widehat D
  \longleftrightarrow
  \ii \widehat H_{\mathrm{cyl}}
  \label{eq:dilatation-cylinder-hamiltonian}
\end{equation}
in Euclidean conventions, or equivalently the eigenvalue of
\(\widehat H_{\mathrm{cyl}}\) is the scaling dimension.

If \(\mathcal O\) has scaling dimension \(\Delta\), then
\[
  \widehat H_{\mathrm{cyl}}\ket{\mathcal O}
  =
  \Delta\ket{\mathcal O}.
\]
Thus the spectrum of scaling dimensions in flat space is the energy spectrum
of the theory on the spatial sphere.  Unitarity implies \(\Delta\ge0\) for
states above the vacuum in a unitary CFT.

\section{Primary and Descendant States}

A local operator at the origin is called a conformal primary if the associated
state is annihilated by all special conformal generators:
\begin{equation}
  \widehat K_\mu\ket{\mathcal O}=0 .
  \label{eq:primary-state-condition}
\end{equation}
If \(\mathcal O_\alpha\) transforms in a finite-dimensional spin
representation with matrices \((S_{\mu\nu})_\alpha{}^\beta\), then a primary
state obeys
\[
  \widehat D\ket{\mathcal O_\alpha}
  =
  \ii\Delta\ket{\mathcal O_\alpha},
  \qquad
  \widehat J_{\mu\nu}\ket{\mathcal O_\alpha}
  =
  -(S_{\mu\nu})_\alpha{}^\beta\ket{\mathcal O_\beta},
  \qquad
  \widehat K_\mu\ket{\mathcal O_\alpha}=0 .
\]
Translations generate descendants:
\[
  \widehat P_{\mu_1}\cdots\widehat P_{\mu_n}\ket{\mathcal O}
  \longleftrightarrow
  \ii^n
  \partial_{\mu_1}\cdots\partial_{\mu_n}\mathcal O(0).
\]
The commutator
\[
  [\widehat D,\widehat P_\mu]=\ii\widehat P_\mu
\]
shows that each derivative raises the scaling dimension by one.

\begin{figure}[H]
  \centering
  \resizebox{0.82\textwidth}{!}{%
  \begin{tikzpicture}[x=1cm,y=1cm,every node/.style={font=\scriptsize}]
    \node[font=\small] at (0,2.65) {conformal multiplet};
    \node[draw=black,rounded corners,inner sep=7pt] (p) at (0,1.75)
      {primary \(\ket{\mathcal O}\)};
    \node[draw=black,rounded corners,inner sep=7pt] (d1) at (-2.4,0.55)
      {\(\widehat P_\mu\ket{\mathcal O}\)};
    \node[draw=black,rounded corners,inner sep=7pt] (d2) at (2.4,0.55)
      {\(\widehat P_\nu\ket{\mathcal O}\)};
    \node[draw=black,rounded corners,inner sep=7pt] (d3) at (0,-0.75)
      {\(\widehat P_\mu\widehat P_\nu\ket{\mathcal O}\)};
    \draw[-{Latex[length=2mm]},thick] (p)--(d1);
    \draw[-{Latex[length=2mm]},thick] (p)--(d2);
    \draw[-{Latex[length=2mm]},thick] (d1)--(d3);
    \draw[-{Latex[length=2mm]},thick] (d2)--(d3);
    \node[blue!65!black] at (-3.8,1.15) {\(+1\) in \(\Delta\)};
    \node[orange!85!black] at (3.85,1.15) {\(\widehat K_\mu\) lowers};
    \node[orange!85!black,align=center] at (0,-1.65)
      {the primary is the lowest-energy state\\in the multiplet};
  \end{tikzpicture}
  }
  \caption{A conformal multiplet is generated from a primary by translations.
  On the cylinder, descendants have energies shifted upward by integers.}
  \label{fig:primary-descendant-multiplet}
\end{figure}

\section{Reflection Positivity and Conjugation}

Radial quantization gives a natural conjugation.  In Euclidean flat space,
reflection through the unit sphere is inversion:
\[
  x^\mu\mapsto \frac{x^\mu}{x^2}.
\]
On the cylinder this is time reflection \(\tau\mapsto-\tau\).  For a scalar
operator of dimension \(\Delta\), the conjugate insertion at infinity is
\[
  \mathcal O^\dagger(\infty)
  =
  \lim_{r\to\infty}r^{2\Delta}\mathcal O(r n)
\]
when the limit is independent of \(n\).  The two-point function becomes the
inner product
\[
  \braket{\mathcal O}{\mathcal O}
  =
  \left\langle
  \mathcal O^\dagger(\infty)\mathcal O(0)
  \right\rangle .
\]
For unitary theories this inner product is positive semidefinite and becomes
positive definite after null states are quotiented.  This positivity is the
source of unitarity bounds, developed after primary transformation laws and
two-point functions have been fixed.

\section{The Free Scalar Check}

For the conformally coupled free scalar on
\(\R\times S^{D-1}\), the scalar curvature of the unit cylinder is
\[
  R=(D-1)(D-2).
\]
With \(\xi_\ast=(D-2)/(4(D-1))\), the conformal mass term is
\[
  m_{\mathrm{cyl}}^2=\xi_\ast R=\frac{(D-2)^2}{4}.
\]
For a one-particle mode on \(S^{D-1}\) with angular momentum \(\ell\),
\[
  -\nabla^2_{S^{D-1}}Y_{\ell}=\ell(\ell+D-2)Y_{\ell}.
\]
The cylinder energy is
\[
  E_\ell
  =
  \sqrt{\ell(\ell+D-2)+\frac{(D-2)^2}{4}}
  =
  \ell+\frac{D-2}{2}.
\]
The \(\ell=0\) mode corresponds to the primary field \(\phi\), whose scaling
dimension is
\[
  \Delta_\phi=\frac{D-2}{2}.
\]
Modes with \(\ell>0\) correspond to derivative descendants of \(\phi\).  This
explicit spectrum realizes the state-operator correspondence in a theory where
all local operators can be built from normal-ordered products of \(\phi\) and
its derivatives, modulo the equation of motion.
